% !TEX TS-program = pdflatex
% !TEX encoding = UTF-8 Unicode

% This is a simple template for a LaTeX document using the "article" class.
% See "book", "report", "letter" for other types of document.

\documentclass[11pt]{article} % use larger type; default would be 10pt

\usepackage[utf8]{inputenc} % set input encoding (not needed with XeLaTeX)

%%% Examples of Article customizations
% These packages are optional, depending whether you want the features they provide.
% See the LaTeX Companion or other references for full information.

%%% PAGE DIMENSIONS
\usepackage{geometry} % to change the page dimensions
\geometry{a4paper} % or letterpaper (US) or a5paper or....
% \geometry{margin=2in} % for example, change the margins to 2 inches all round
% \geometry{landscape} % set up the page for landscape
%   read geometry.pdf for detailed page layout information
\usepackage{bm}
\usepackage{graphicx} % support the \includegraphics command and options
\usepackage{tabularx}  % in the preamble
% \usepackage[parfill]{parskip} % Activate to begin paragraphs with an empty line rather than an indent

%%% PACKAGES
\usepackage{booktabs} % for much better looking tables
\usepackage{array} % for better arrays (eg matrices) in maths
\usepackage{paralist} % very flexible & customisable lists (eg. enumerate/itemize, etc.)
\usepackage{verbatim} % adds environment for commenting out blocks of text & for better verbatim
\usepackage{subfig} % make it possible to include more than one captioned figure/table in a single float
% These packages are all incorporated in the memoir class to one degree or another...

%%% HEADERS & FOOTERS
\usepackage{fancyhdr} % This should be set AFTER setting up the page geometry
\pagestyle{fancy} % options: empty , plain , fancy
\renewcommand{\headrulewidth}{0pt} % customise the layout...
\lhead{}\chead{}\rhead{}
\lfoot{}\cfoot{\thepage}\rfoot{}

%%% SECTION TITLE APPEARANCE
\usepackage{sectsty}
\allsectionsfont{\sffamily\mdseries\upshape} % (See the fntguide.pdf for font help)
% (This matches ConTeXt defaults)
\usepackage{rotating} % For rotating tables
%%% ToC (table of contents) APPEARANCE
\usepackage[nottoc,notlof,notlot]{tocbibind} % Put the bibliography in the ToC
\usepackage[titles,subfigure]{tocloft} % Alter the style of the Table of Contents
\renewcommand{\cftsecfont}{\rmfamily\mdseries\upshape}
\renewcommand{\cftsecpagefont}{\rmfamily\mdseries\upshape} % No bold!
\usepackage[lf]{Minion Pro}
\usepackage{Myriad Pro}



% definitions for citeproc citations
\NewDocumentCommand\citeproctext{}{}
\NewDocumentCommand\citeproc{mm}{%
	\begingroup\def\citeproctext{#2}\cite{#1}\endgroup}
\makeatletter
% allow citations to break across lines
\let\@cite@ofmt\@firstofone
% avoid brackets around text for \cite:
\def\@biblabel#1{}
\def\@cite#1#2{{#1\if@tempswa , #2\fi}}
\makeatother
\newlength{\cslhangindent}
\setlength{\cslhangindent}{1.5em}
\newlength{\csllabelwidth}
\setlength{\csllabelwidth}{3em}
\newenvironment{CSLReferences}[2] % #1 hanging-indent, #2 entry-spacing
{\begin{list}{}{%
			\setlength{\itemindent}{0pt}
			\setlength{\leftmargin}{0pt}
			\setlength{\parsep}{0pt}
			% turn on hanging indent if param 1 is 1
			\ifodd #1
			\setlength{\leftmargin}{\cslhangindent}
			\setlength{\itemindent}{-1\cslhangindent}
			\fi
			% set entry spacing
			\setlength{\itemsep}{#2\baselineskip}}}
	{\end{list}}

%%% END Article customizations

%%% The "real" document content comes below...

\title{Brief Article}
\author{The Author}
%\date{} % Activate to display a given date or no date (if empty),
         % otherwise the current date is printed 



\begin{document}

\begin{table}
\begin{tabular}{|l|c|c|c|c|} \hline
 \bfseries Participant & $\bm{F_{CE}^{max}}$ [N] & $\bm{L_{CE}^{opt}}$ [cm] & $\bm{L_{SEE}^{0}}$ [cm] & \bfseries RMSD [\%] \\ \hline
1 & 4400 & 8.6 & 17 & 7.3 \\ \hline
2 & 5100 & 9.9 & 15 & 6.9 \\ \hline
3 & 5600 & 8.7 & 17 & 5.8 \\ \hline
4 & 6800 & 8.6 & 16 & 4.1 \\ \hline
5 & 9200 & 7.4 & 18 & 4.7 \\ \hline
6 & 8500 & 8.8 & 16 & 2.9 \\ \hline
AVG±SD & 6600 $\pm$ 1700 & 8.7 $\pm$ 0.7 & 17 $\pm$ 1 & 5.3 $\pm$ 1.6 \\ \hline
\end{tabular}\break\hfill\footnotesize{$F_{CE}^{max}$, $L_{CE}^{opt}$ and $L_{SEE}^{0}$ were estimated from the isometric measurements for each participant individually. The root mean square difference (RMSD) was normalised by the individual maximal predicted knee joint moment, which is the product of $F_{CE}^{max}$ and the moment arm of \textit{m. quadriceps femoris.}}
\end{table}

\begin{table}
 \begin{tabularx}{\textwidth}{|p{1.0cm}|*{9}{>{\centering\arraybackslash}X|}} \hline
 & \multicolumn{9}{c|}{\textit{Condition}} \\ \hline
 & \bfseries 0.20 & \bfseries 0.25 & \bfseries 0.30 & \bfseries 0.35C & \bfseries 0.40 & \bfseries 0.35A & \bfseries 0.35B & \bfseries 0.35D & \bfseries 0.35E \\ \hline
 \bm{$W_{ext}^{+}$}  & 102.3 $\pm$ 2.5  & 102.2 $\pm$ 2.0  & 103.4 $\pm$ 3.2  & 103.3 $\pm$ 4.3  & 104.2 $\pm$ 3.5  & 104.5 $\pm$ 4.0  & 102.5 $\pm$ 3.6  & 101.9 $\pm$ 4.2  & 100.5 $\pm$ 2.7 \\ \hline 
 \bm{$W_{flex}^{+}$}  & 1.3 $\pm$ 1.2  & 0.9 $\pm$ 0.9  & 1.1 $\pm$ 1.2  & 0.9 $\pm$ 0.7  & 1.4 $\pm$ 1.0  & 1.6 $\pm$ 1.5  & 1.7 $\pm$ 1.2  & 2.6 $\pm$ 1.9  & 3.3 $\pm$ 1.8 \\ \hline 
 \bm{$W_{ext}^{-}$}  & -0.0 $\pm$ 0.1  & -0.0 $\pm$ 0.0  & -0.0 $\pm$ 0.0  & -0.0 $\pm$ 0.0  & -0.0 $\pm$ 0.0  & -0.0 $\pm$ 0.0  & -0.0 $\pm$ 0.0  & -0.0 $\pm$ 0.0  & -0.1 $\pm$ 0.4 \\ \hline 
 \bm{$W_{flex}^{-}$}  & -3.7 $\pm$ 2.2  & -3.1 $\pm$ 2.0  & -4.6 $\pm$ 3.3  & -4.3 $\pm$ 4.1  & -5.6 $\pm$ 3.4  & -6.0 $\pm$ 3.6  & -4.3 $\pm$ 3.0  & -4.7 $\pm$ 4.4  & -3.8 $\pm$ 2.0 \\ \hline 
\end{tabularx}\break\hfill\footnotesize{Positive (+) and negative (-) mechanical work during knee extension (ext) and flexion (flex) were calculated by integrating the instantaneous mechanical power over time, separately for the periods when the instantaneous power was positive (positive work) and negative (negative work) respectively. }
\end{table}

\begin{table}
 \begin{tabularx}{\textwidth}{|p{2.8cm}|*{9}{>{\centering\arraybackslash}X|}} \hline
  & \multicolumn{9}{|c|}{\itshape Condition} \\ \hline
  & \bfseries 0.20 & \bfseries 0.25 & \bfseries 0.30 & \bfseries 0.35C & \bfseries 0.40 & \bfseries 0.35A & \bfseries 0.35B & \bfseries 0.35D & \bfseries 0.35E \\ \hline
 \multicolumn{10}{|r|}{\itshape Movement characteristics} \\ \hline
 \bfseries CF [Hz]  & 0.44  & 0.59  & 0.75  & 0.93  & 1.13  & 0.79  & 0.83  & 1.13  & 1.35 \\ \hline 
 \bfseries FTS [ ]  & 0.35  & 0.40  & 0.45  & 0.51  & 0.58  & 0.72  & 0.61  & 0.46  & 0.48 \\ \hline 
 \bm{$T_{ext}$} \bfseries [s]  & 0.79  & 0.67  & 0.60  & 0.55  & 0.51  & 0.91  & 0.73  & 0.41  & 0.36 \\ \hline 
 \bm{$T_{flex}$} \bfseries [s]  & 1.46  & 1.02  & 0.73  & 0.53  & 0.38  & 0.36  & 0.48  & 0.48  & 0.38 \\ \hline 
 \multicolumn{10}{|r|}{\itshape AMPO} \\ \hline
 \bfseries Participant 1 [s\textsuperscript{-1}]  & 0.17  & 0.23  & 0.30  & 0.36  & 0.44  & 0.34  & 0.36  & 0.39  & 0.40 \\ \hline 
 \bfseries Participant 2 [s\textsuperscript{-1}]  & 0.17  & 0.23  & 0.26  & 0.32  & 0.37  & 0.32  & 0.30  & 0.32  & 0.35 \\ \hline 
 \bfseries Participant 3 [s\textsuperscript{-1}]  & 0.16  & 0.19  & 0.23  & 0.29  & 0.36  & 0.27  & 0.29  & 0.32  & 0.34 \\ \hline 
 \bfseries Participant 4 [s\textsuperscript{-1}]  & 0.17  & 0.21  & 0.27  & 0.30  & 0.35  & 0.30  & 0.29  & 0.30  & 0.35 \\ \hline 
 \bfseries Participant 5 [s\textsuperscript{-1}]  & 0.14  & 0.17  & 0.22  & 0.25  & 0.30  & 0.26  & 0.26  & 0.26  & 0.26 \\ \hline 
 \bfseries Participant 6 [s\textsuperscript{-1}]  & 0.17  & 0.20  & 0.24  & 0.29  & 0.33  & 0.31  & 0.30  & 0.30  & 0.32 \\ \hline 
 \bfseries Average $\pm$ SD [s\textsuperscript{-1}]  & 0.16 ± 0.01  & 0.21 ± 0.02  & 0.25 ± 0.03  & 0.30 ± 0.03  & 0.36 ± 0.04  & 0.30 ± 0.03  & 0.30 ± 0.03  & 0.32 ± 0.04  & 0.34 ± 0.04 \\ \hline 
 \bfseries Predicted [s\textsuperscript{-1}]  & 0.20  & 0.25  & 0.30  & 0.35  & 0.40  & 0.35  & 0.35  & 0.35  & 0.35 \\ \hline 
\end{tabularx}\break\hfill\footnotesize{CF denotes cycle frequency, FTS denotes fraction of cycle time spent shortening, $T_{ext}$ and $T_{flex}$ denote the knee joint extension and flexion time per cycle. Measured AMPO was normalised to the product of the participant's individually estimated $F_{CE}^{max}$ and the generic model parameter value of $L_{CE}^{opt}$, while predicted AMPO was normalised to the product of model parameters $F_{CE}^{max}$ and $L_{CE}^{opt}$ (see Methods). }
\end{table}

\begin{table}
	\begin{tabular}{|l|l|l|r|l|} \hline
  \bfseries Description & \bfseries Symbol & \bfseries Unit & \bfseries Value & \bfseries Reference \\ \hline
  \multicolumn{5}{|r|}{\itshape Musculoskeletal geometry} \\ \hline
MTC length at $\phi_{knee}=0$ & $A_0$ & cm & 21.3 & \citeproc{ref-van_soest_which_2000}{van Soest \& Casius (2000)} \\ \hline
Linear coefficient of MTC-$\phi_{knee}$ relation & $A_1$ & cm/rad & 4.20 & \citeproc{ref-van_soest_which_2000}{van Soest \& Casius (2000)} \\ \hline
Hill constant & $a^{rel}$ & - & 0.41 & \citeproc{ref-van_soest_which_2000}{van Soest \& Casius (2000)} \\ \hline
  \multicolumn{5}{|r|}{\itshape Contraction dynamics} \\ \hline
Hill constant & $b^{rel}$ & 1/s & 5.20 & \citeproc{ref-van_soest_which_2000}{van Soest \& Casius (2000)} \\ \hline
Maximum isometric CE force & $F_{CE}^{max}$ & N & 5,250 & \citeproc{ref-van_soest_which_2000}{van Soest \& Casius (2000)} \\ \hline
PEE stiffness scaling parameter & $k_{PEE}$ & N/m<sup>2</sup> & 1.19 $\times$ 10\textsuperscript{7} & \citeproc{ref-van_soest_contribution_1993}{van Soest \& Bobbert (1993)} \\ \hline
SEE stiffness scaling parameter & $k_{SEE}$ & N/m<sup>2</sup> & 1.28 $\times$ 10\textsuperscript{8} & \citeproc{ref-van_soest_which_2000}{van Soest \& Casius (2000)} \\ \hline
CE optimum length & $L_{CE}^{opt}$ & cm & 9.30 & \citeproc{ref-van_soest_which_2000}{van Soest \& Casius (2000)} \\ \hline
PEE slack length & $L_{PEE}^0$ & cm & 13.0 & \citeproc{ref-van_soest_contribution_1993}{van Soest \& Bobbert (1993)} \\ \hline
SEE slack length & $L_{SEE}^0$ & cm & 16.0 & \citeproc{ref-van_soest_which_2000}{van Soest \& Casius (2000)} \\ \hline
  \multicolumn{5}{|r|}{\itshape Excitation dynamics} \\ \hline
Determines CE force-length relation width & $w$ & - & 0.56 & \citeproc{ref-van_soest_which_2000}{van Soest \& Casius (2000)} \\ \hline
Determines steepness $\gamma$-$q$ relation & $a_{act}$ & - & -4.59 & \citeproc{ref-hatze_myocybernetic_1981}{Hatze (1981)} \\ \hline
Determines $pCa^{2+}$ level at which $q=0.5$ & $b_{act,1}$ & - & 5.17 & \citeproc{ref-hatze_myocybernetic_1981}{Hatze (1981)} \\ \hline
Determines $pCa^{2+}$ level at which $q=0.5$ & $b_{act,2}$ & - & 1.08 & \citeproc{ref-hatze_myocybernetic_1981}{Hatze (1981)} \\ \hline
Determines $pCa^{2+}$ level at which $q=0.5$ & $b_{act,3}$ & - & -0.19 & \citeproc{ref-hatze_myocybernetic_1981}{Hatze (1981)} \\ \hline
Relates $\gamma$ to $pCa^{2+}$ & $kCa$ & mol/L & 8.00 $\times$ 10\textsuperscript{-6} & \citeproc{ref-kistemaker_length-dependent_2005}{Kistemaker et al. (2005)} \\ \hline
Activation time constant & $\tau_{act}$ & ms & 88.9 & \citeproc{ref-hatze_myocybernetic_1981}{Hatze (1981)} \\ \hline
De-activation time constant & $\tau_{deact}$ & ms & 88.9 & \citeproc{ref-hatze_myocybernetic_1981}{Hatze (1981)} \\ \hline
\end{tabular}\break\hfill\footnotesize{Parameter values are depicted only for those key to predict the maximally attainable AMPO. }
\end{table}




\begin{table}[htbp]
\centering
\caption{Your table caption}
\begin{tabularx}{\textwidth}{|l|X|r|r}
\hline
Label & Description & 0.40 & 0.35C  \\
\hline
WposExt & Work positive extension & 12.1 ± 1.2 & 11.5 ± 1.1  \\
WposFlx & Work positive flexion & 8.3 ± 0.9 & 7.9 ± 0.8 . \\
\hline
\end{tabularx}
\end{table}
    
\begin{table}[htbp]
\centering
\caption{Auto-stretching table}
 \begin{tabularx}{\textwidth}{|l|*{9}{>{\centering\arraybackslash}X|}} \hline
Label & 0.40 & 0.35C & 0.30 & 0.25 &  0.20 & 0.35A & 0.35B & 0.35D & 0.35E \\ \hline
Label & 0.40 & 0.35C & 0.30 & 0.25 & 0.20 & 0.35A & 0.35B & 0.35D & 0.35 \\
\hline
WposExt & 12.1 ± 1.2 & 13.2 ± 1.1 & 13.2 ± 1.1 & 13.2 ± 1.1 & 13.2 ± 1.1 & 13.2 ± 1.1 & 13.2 ± 1.1 & 13.2 ± 1.1 & 13.2 ± 1.1 \\
% other rows
\hline
\end{tabularx}
\end{table}



\end{document}
